\section{Relação com trabalhos anteriores}
\label{sec:resultado-literatura}

Esta seção compara os resultados com as ferramentas e os estudos do Capítulo~\ref{cap:trabalhos-relacionados}. A comparação tem um limite de origem: aqui há opinião e preferência declarada, enquanto parte dos trabalhos citados mediu desempenho.

A preferência dos estudantes pela versão em português e a concordância do painel com a afirmação de que palavras em português facilitam a compreensão, com média 4,38, vão na direção da escolha feita pelo Portugol Studio, que traduz as palavras-chave para que o iniciante não precise lidar com dois idiomas ao mesmo tempo \cite{univali2026}. Os dados não alcançam uma diferença de desenho entre as duas propostas. No Portugol o vocabulário em português é fixo; na plataforma ele é escolhido por quem usa. O formulário dos estudantes comparou uma versão em inglês com uma em português definida pelos autores e, por isso, não informa o efeito de o próprio aluno escolher as palavras.

\citeonline{stefik2013} compararam o desempenho de novatos em linguagens consolidadas com o obtido em uma linguagem de controle de palavras-chave aleatórias e encontraram efeito da forma sintática sobre o desempenho. Os dados desta avaliação são de outro tipo, porque não houve tarefa nem contagem de erros. Eles não confirmam nem contrariam aquele resultado, e a comparação depende do estudo de desempenho previsto no Capítulo~\ref{cap:consideracoes}.

Cinco dos oito especialistas conheciam o Portugol Studio ou o Portugol Webstudio, e mesmo assim a afirmação de que a plataforma oferece algo que as ferramentas conhecidas não oferecem teve média 4,62. O formulário não pergunta qual seria esse diferencial. A execução passo a passo, item de maior valor didático no painel, com média 4,75, também existe no Portugol Studio, que tem depurador e inspetor de variáveis \cite{univali2026}. O que a plataforma acrescenta é a configuração do vocabulário pelo assistente e a exibição das fases da tradução. O assistente foi bem avaliado, mas um respondente teve dificuldade para encontrar o botão de criar linguagem, e as fases da tradução foram pouco visitadas: dois especialistas declararam ter visto a tabela de \textit{tokens} e quatro, o código intermediário.

A correção automática pelos casos de teste, com média 4,50 entre os que lecionam, segue o ciclo de submissão e veredito de juízes \textit{online} como o Beecrowd \cite{beecrowd2026faq}. A diferença é que os casos são executados com a configuração de linguagem do aluno, como descrito na Seção~\ref{sec:sistema-correcao}. Por fim, a preocupação com a passagem para ferramentas de mercado, que dividiu o painel no item de migração, também orienta o Portugol Studio, cuja tela reproduz a organização de uma IDE profissional para que a familiaridade seja transferível \cite{univali2026}. Lá a preocupação recai sobre a interface; aqui o painel a levantou sobre a sintaxe.

\section{Limites de interpretação e continuidade}
\label{subsec:limites-interpretacao}

Além do limite já citado sobre o público, duas ressalvas afetam a leitura dos resultados do painel. A primeira é sobre a redação em português dos itens do SUS: este trabalho não registra a conferência dela com uma tradução validada, o que pede cautela ao comparar o escore com a média histórica. A segunda é sobre a análise: escores altos dados por especialistas não comprovam aprendizagem, redução de carga cognitiva nem aumento de engajamento.

A sessão com estudantes tem limites próprios. Foi uma única turma, escolhida por conveniência. A opinião foi colhida logo depois de uma demonstração feita pelos próprios autores e de um período de uso da plataforma, situação em que a novidade e a tendência a concordar com quem apresenta pesam nas respostas. Quase todos os alunos já tinham programado, e o formulário não mede desempenho. As respostas indicam aceitação e preferência declarada. Não mostram que a barreira sintática diminuiu nem que houve aprendizagem.

Investigar a hipótese pedagógica é a etapa seguinte e exige outro desenho de estudo: comparar, na mesma IDE, o vocabulário padrão e o vocabulário personalizado, com tarefas equivalentes de entrada e saída, seleção e repetição. A ordem das condições deve ser alternada entre os participantes para reduzir o efeito de prática, e os enunciados, os casos de teste, o tempo disponível e os critérios de classificação dos erros precisam ser definidos antes da coleta, separando erros de escrita de falhas de lógica e de dificuldades de navegação. Com isso seria possível descrever os acertos, o tempo até a primeira execução correta e a frequência de erros em cada condição. Ainda assim, uma diferença nessas medidas não comprova, sozinha, que houve aprendizagem duradoura. Esse estudo deve ter aprovação do comitê de ética antes de qualquer coleta.

O formulário aplicado não tem itens sobre acessibilidade, e as barreiras registradas na Seção~\ref{sec:acessibilidade-responsividade} continuam entre as pendências deste trabalho. Antes de uma nova coleta com participantes, é preciso corrigi-las e ampliar os testes de teclado para os fluxos de salvar linguagem, editar código, usar o terminal e enviar exercícios. A verificação deve cobrir foco visível, entrada e saída de diálogos, nomes dos controles, anúncio de erros e compatibilidade com leitor de tela, registrando qual sistema operacional, navegador e tecnologia assistiva foram usados. Também devem entrar os temas claro e escuro e as larguras de 320, 390, 768 e 1440 pixels CSS. Os critérios WCAG apresentados no referencial orientam essa inspeção, mas declarar conformidade exigiria examinar as páginas e os processos completos \cite{w3cwcag22,w3capg2026}.

\section{Síntese e confronto com os objetivos}
\label{sec:resultado-sintese}

Esta seção retoma os objetivos específicos da introdução e diz, para cada um, o que a avaliação sustenta e o que continua sem verificação.

O primeiro objetivo era especificar os elementos configuráveis da linguagem. O painel considerou o conjunto adequado ao propósito da ferramenta, com média 4,57 entre sete respondentes e nenhuma nota abaixo de 4. A pergunta sobre construções que fizeram falta teve sete respondentes, porque R4 a deixou em branco. Cinco marcaram não ter sentido falta de nada relevante para o ensino introdutório, dois marcaram registros ou estruturas compostas, e modularização em vários arquivos, orientação a objetos e manipulação de cadeias mais rica tiveram uma marcação cada. R7 marcou ao mesmo tempo a primeira alternativa, registros e orientação a objetos, e a resposta foi mantida como veio. Entre os que lecionam, dois dos três que responderam consideraram o subconjunto suficiente para uma disciplina introdutória. A especificação não foi contestada, e as ausências apontadas coincidem com as limitações já declaradas no Capítulo~\ref{cap:consideracoes}.

O segundo objetivo era projetar e implementar a IDE com o assistente, a edição, a execução e o retorno ao usuário. O escore SUS de 75,6, indicativo, e as médias entre 4,00 e 4,57 nos itens do assistente sustentam que o painel conseguiu usar a interface e não contestou as decisões de dividir a personalização em etapas, indicar a etapa corrente e mostrar a pré-visualização. Dois achados vão no outro sentido. O acesso ao assistente falhou para pelo menos um respondente, que não encontrou com facilidade o botão de criar linguagem, e um professor terminou a visita sem ter aberto o assistente. Além disso, na dimensão pragmática do UEQ-S, os pares complicada ou fácil e confusa ou clara tiveram as menores médias, e o item do SUS ``achei a plataforma fácil de usar'' ficou em 3,75. Metade do painel usou a plataforma por menos de 15 minutos, e isso vale tanto para as notas altas quanto para as baixas.

O terceiro objetivo era integrar a IDE ao interpretador. Os respondentes escreveram e executaram programas, viram o código intermediário e usaram a execução passo a passo, e avaliaram a ferramenta como rápida e estável, com média 4,88. A afirmação de que o código intermediário exibido corresponde ao que o programa executa teve média 4,83 entre seis respondentes. Trata-se de percepção de uso. A verificação da integração está nos testes descritos na Seção~\ref{sec:testes-automatizados}.

O quarto objetivo era disponibilizar a plataforma em ambiente \textit{web}, com os fluxos de uso concluídos e com atenção a acessibilidade e responsividade. A avaliação mostra que respondentes externos ao projeto percorreram os fluxos de personalização, execução e submissão. Os fluxos de turma foram pouco exercitados: dois respondentes declararam ter entrado em uma turma ou criado uma. O formulário não tem itens de acessibilidade, e esse requisito continua no estado descrito na Seção~\ref{sec:acessibilidade-responsividade}.

O quinto objetivo era propor um protocolo de avaliação. O protocolo foi proposto e aplicado, com oito respondentes, e a ele se somou a sessão com dezoito estudantes. A aplicação mostrou limites do próprio instrumento: os itens do bloco de uso em sala foram respondidos por quem pouco usou os recursos de turma.

Quanto à hipótese pedagógica, vale a distinção da Seção~\ref{subsec:limites-interpretacao}. O que foi coletado é opinião. A maioria do painel considera plausível que a personalização reduza a barreira sintática, com média 3,88, e um respondente discorda das três premissas. Onze de dezoito estudantes acharam a versão em português mais fácil de entender, depois de uma demonstração e de um período de uso, sem tarefa medida. Nenhum desses dados prova a redução da barreira nem aprendizagem. Do lado desfavorável, a intenção de adotar a plataforma em uma disciplina introdutória ficou em 3,00 entre os quatro que lecionam, um professor não vê espaço para ela na disciplina, o painel se dividiu sobre o risco de migração para uma linguagem de mercado e três de sete respondentes concordam que a personalização pode desviar a atenção do raciocínio algorítmico. Esses pontos têm o mesmo peso dos resultados favoráveis e orientam os encaminhamentos do Capítulo~\ref{cap:consideracoes}.
